\chapter{Measure Theory}

\section{Measurable Spaces}
2
\subsection{Problems}
\begin{problem}
	Partition generated by a $\sigma$-algebras.
	\begin{enumerate}
		\item Let $\mathcal{C} = \{A,B,C\}$ be a partition of $E$ List all the elements of $\sigma \mathcal{C}$
		\item let $\mathcal{C}$ be a countable partition of $E$. show that every element of $\sigma \mathcal{C}$ is a countable union of elements taken from $\mathcal{C}$.
		\item Let $E = \mathbb{R}$. Let $\mathcal{C}$ be the collection of all singleton subsets of $\mathbb{R}$, that is, sets that only contain a single point. Show that each element of $\sigma \mathcal{C}$ is either countable, or the complement of a countable set.  
	\end{enumerate}
\end{problem}

\begin{solution}
	% (a)
	
	\textbf{1.} The elements of $\sigma \mathcal{C}$ are defined as the intersection of all sigma-algebras that contain $\mathcal{C}$. Since $\mathcal{C}$ contains 3 subsets of $E$, as it is a partition, which implies that $A,B,C$ are disjoint and have a union of $E$. Thus, we can define $\sigma \mathcal{C}$ as the following set:
	\begin{equation}
		\sigma \mathcal{\mathcal} = \{\emptyset, E, \{A\}, \{B\}, \{C\}, \{A\cap B\}, \{A\cap C\}, \{B\cap C\}\big\}
	\end{equation}
	This set contains $E$, all partitions of closures by $E$ generated from $\mathcal{C}$. Thus, this is the smallest sigma-algebra that contains $E$ generated by $\mathcal{C}$.
	
	\textbf{2.} Since $\mathcal{C}$ is a partition,  $\bigcup \mathcal{C} = E$. We need to show that $\sigma \mathcal{C}$ contains only elements from $C$. We can take $\varepsilon$ as the set of all countable unions taken from $E$. We can show that this is a sigma-algebra by taking its definitions. Since $\varepsilon$ contains $E$ the subset of all the partitions, it must also contain the union of none of the partitions, or $\emptyset$. Secondly, by taking all of the unions, we start with the disjoint partitions and then create their unions until we reach $E$. This will be closed under complement as by taking countable unions, we will have the sets that are not part of the union as another subset of a different union. Thus, we can argue that this is the same as $\sigma \mathcal{C}$ by the logic of construction. Thus, since each element of  $\sigma \mathcal{C}$ is isomorphic to the construction that we obtained from the union of subsets, we have shown that all elements are just a set of unions in $\mathcal{C}$.
	
	\textbf{3.} Define the set $\mathcal{A} = \{ A \subset \mathbb{R} : A \text{ is countable or } A^c \text{ is countable}\}$. We want to show that this is a sigma-algebra. Since $\emptyset$ is countable, it is contained within $\mathcal{A}$. $\emptyset \in \mathcal{A}$. The complement of this is the set of reals; $\mathbb{R} = \emptyset^c \in \mathcal{A}$. By definition, if $A$ is countable, then $A^c$ is co-countable. Similarly, if $A$ is co-countable, this implies that $A^c$ is countable. Thus, both sets are contained in $\mathcal{A}$. 
	
	Thus, consider a sequence $(A_n)_{n\in\mathbb{N}} \subset \mathcal{A}$. the union of this sequence is $\bigcup_n A_n$. This is split into cases. First, if all $A_n$ are countable, so it the union, which is included in $A_n$. If all $A_n$ are co-countable, then consider the compliment of the union $(\bigcup_n A_n)^c = \bigcap_n A_n^c$ which is an intersection of countable sets, which is countable. If there is a mix between countable and co-countable sets, they are co-countable, which are contained in $A$ by the previous definition. Thus, $\bigcup_n A_n \in \mathcal{A}$. Thus, $\mathcal{A}$ is closed under union and complement. Thus, $\mathcal{A}$ is a sigma-algebra. Thus, $\sigma \mathcal{C} \subset \mathcal{A}$
	
	We want to show that any countable set is a union of the singletons. Since $A$ contains countable and co-countable sets, $A$ is a union of singletons, which are by definition contained in $\mathcal{A}$.
	\begin{equation}
		A = \bigcup_{x\in A}\{x\}
	\end{equation}
	For co-countable sets in $A$, we note that the complements are countable, which are contained in the above definition. Since $\mathcal{A}$ consists of countable sets, these are the countable unions of singletons contained in $\sigma \mathcal{C}$. Since $\mathcal{A}$ also consists of co-countable sets, these are the complements of intersections of countable sets in $\sigma \mathcal{C}$. Thus, $\mathcal{A} \subset \sigma \mathcal{C}$.
\end{solution}